\documentclass[12pt]{article}
\usepackage[utf8]{inputenc}
\usepackage[T1]{fontenc}
\usepackage{hyperref}
\title{Dacian Language Resources for K-12 Curriculum Development}
\author{Compiled by Claude Sonnet 5.5 (claude-sonnet-5-5) for Andr\'as Kornai}
\date{}
\begin{document}
\maketitle

\section{Introduction}
Dacian is an extinct Indo-European language (ISO 639-3 xdc; Glottocode daci1234) of ancient Dacia. It has no native speakers and no continuous tradition. Search results show it is very poorly documented. Only one inscription is believed to survive. The evidence consists mostly of place names, river names, personal names and some plant names. Any curriculum must therefore be a heritage and awareness curriculum. It cannot teach Dacian as a language of instruction.

\section{Primary Educational Resources}
I found no primers, textbooks or school programmes for Dacian. No Dacian-English dictionary or word list turned up in the searches, other than the Wikipedia pages listed below.

\section{Linguistic and Reference Materials}
\begin{itemize}
\item Wikipedia, ``Dacian language''. It summarizes the attestation: one inscription, about 60 plant names in Dioscorides, and about 100 words reconstructed by comparative linguistics. \url{https://en.wikipedia.org/wiki/Dacian_language}
\item Wikipedia, ``List of Romanian words of possible pre-Roman origin''. It cautions that it is difficult to verify whether a given Romanian word is actually from Dacian. \url{https://en.wikipedia.org/wiki/List_of_Romanian_words_of_possible_pre-Roman_origin}
\item A Springer chapter that describes Dacian as a mystery, with direct evidence limited mostly to toponyms, hydronyms and personal names. \url{https://link.springer.com/chapter/10.1007/978-3-031-04051-1_27}
\item A Romanian journal article on a dictionary and the substratum question, ``Inherited from the substratum, ultimately''. \url{https://journals.unibuc.ro/index.php/tal/article/view/1346}
\end{itemize}

\section{Religious and Community Resources}
None found in the searches.

\section{Digital and Audio Resources}
None found. No audio or digital corpus exists, because the language is not directly attested beyond names and glosses.

\section{Curriculum Development Context}
Little exists for Dacian, so the dossier is short. Teachers could use the reconstructed word lists, the toponyms and the Romanian substratum words listed above. They should state clearly that the forms are reconstructed and uncertain. No community revitalization programme was found.
\end{document}